% s-104-p104-m002; original PDF 104, 105
\Aphorism{s-104-p104-m002}{101}
Dass der Mann sich um die Liebe des Weibes bewerben soll, ist für die Ausnahme-Männer ein hartes Gesetz; es beweist, dass das Weib sich nicht viel Kritik und Spürsinn zutraut (sonst würde es suchen, statt sich suchen zu lassen), und dass mithin gerade der Ausnahmemensch geringe Chancen hat, gewählt zu werden. Aber schon, dass er überhaupt mit Anderen in Wettbewerb treten soll, wird den höheren Menschen schwer ankommen; er würde es natürlicher finden, wenn das Weib ihm von selbst und ohne grossen Zeitverlust zufiele. Seinem Geschmacke würde etwa folgende Weibslogik zusagen: „Du hast andere Siege zu feiern, als den über mich; bin ich dazu gemacht, mich dir vorzuenthalten? Mein Widerstand würde deinen Willen in einer Richtung anstacheln, wo ihm Enttäuschung gewiss ist; ich bin kein Exempel für einen Rechner wie dich. Was hätte ich auch an dich zu verlieren, was nicht schon mein ganzes Geschlecht an dich verloren hätte! Du bist über uns hinaus; deine Ungeduld, die noch ganz andere Probleme zur Verzweiflung getrieben hat, errieth allzuschnell, wie es mit uns bestellt ist. Wir sind zu kurzathmig für deinen Sturmwind und zu sehr Kork für die Tiefe deiner Freundschaft; unser ewig-Weibliches zieht euch hinan, nämlich vom Grunde zur Oberfläche. Wenn du einer von uns deine Liebe schenkst, so ist das ein \Foreign{faute de mieux}, so gut wir es verlangen können. Und dass du gerade zu mir kommst, ist eine feinere Höflichkeit; ich würde wie ein trotziges Kind handeln, wenn ich sie zurückwiese. Freilich bist du vielleicht als Mann weniger Mann als Andere; aber was hilft’s? Ich fürchte dich noch in deiner Unmännlichkeit: vielleicht ist auch sie nur Ironie und Höflichkeit? Eine Andeutung, dass du mich nicht bloss zur Concubine haben willst, dass du mich, wie euer Philosoph Kant sagt, nicht nur als Mittel, sondern zugleich als Zweck „an sich selbst“ behandeln wirst?“ So ungefähr sollte, bei richtiger Selbsteinschätzung, das Weib denken; statt dessen hält es sich, seiner Unentbehrlichkeit wegen, für ein wirkliches Endziel, vielleicht für das begehrenswertheste, des männlichen Willens, und — wie sagt doch Zarathustra? „auch die hohlste Nuss will noch geknackt sein!“
\par

